
\documentclass[12pt,a4paper]{article}

\usepackage{amsmath}
\usepackage{amssymb}
\usepackage{array}
\usepackage{multirow}
\usepackage{graphicx}
\usepackage[table]{xcolor}
\usepackage{colortbl}
\usepackage{booktabs}
\usepackage{tabularx}
\usepackage{geometry}

\geometry{margin=0.7in}

\definecolor{headerblue}{RGB}{40,78,130}
\definecolor{lightblue}{RGB}{220,235,248}
\definecolor{greenmark}{RGB}{210,240,215}
\definecolor{yellowmark}{RGB}{255,245,190}
\definecolor{redmark}{RGB}{250,215,215}

\begin{document}
	
	\begin{center}
		{\LARGE \textbf{Assignment 3}}\\[5pt]
		{\large Mathematics and Tables Using LaTeX}
	\end{center}
	
	\section*{Que 1}
	
	Write LaTeX code to produce the following expression:
	
	\[
	\overbrace{a+b+c}^{\text{First Group}}
	+
	\underbrace{x+y+z}_{\text{Second Group}}
	=
	\overbrace{a+\underbrace{b+c}_{\text{Sum 1}}}^{\text{Total}}
	\]
	
	\section*{Que 2}
	
	Create the following array using the array environment:
	
	\[
	A=
	\left[
	\begin{array}{c|ccc}
		& x_1 & x_2 & x_3\\
		\hline
		y_1 & a_{11} & a_{12} & a_{13}\\
		y_2 & a_{21} & a_{22} & a_{23}\\
		y_3 & a_{31} & a_{32} & a_{33}
	\end{array}
	\right]
	\]
	
	\section*{Que 3}
	
	Write the following piecewise function using the cases environment:
	
	\[
	f(x)=
	\begin{cases}
		x^2+2x+1, & x<0,\\[6pt]
		\sqrt{x+1}, & 0\leq x<4,\\[6pt]
		\dfrac{x^2-1}{x-1}, & x\geq 4.
	\end{cases}
	\]
	
	\section*{Que 4}
	
	Typeset the following matrix and label its different sections using overbrace:
	
	\[
	\begin{bmatrix}
		\overbrace{a_{11}\quad a_{12}\quad a_{13}}^{\text{Column Group 1}}
		&
		\overbrace{a_{14}\quad a_{15}}^{\text{Column Group 2}}
		\\[8pt]
		a_{21}\quad a_{22}\quad a_{23}
		&
		a_{24}\quad a_{25}
		\\[8pt]
		a_{31}\quad a_{32}\quad a_{33}
		&
		a_{34}\quad a_{35}
	\end{bmatrix}
	\]
	
	\section*{Que 5}
	
	Create the following expression:
	
	\[
	\begin{aligned}
		F(x)=
		\underbrace{
			\sum_{i=1}^{n}
			\left(a_i x_i^2+b_i x_i+c_i\right)
		}_{\text{Quadratic Contribution}}
		+
		\underbrace{
			\begin{bmatrix}
				x_1 & x_2 & \cdots & x_n
			\end{bmatrix}
			A
			\begin{bmatrix}
				x_1\\
				x_2\\
				\vdots\\
				x_n
			\end{bmatrix}
		}_{\text{Matrix Contribution}}
	\end{aligned}
	\]
	
	\[
	\underbrace{F(x)}_{\text{Total Energy}}
	\]
	
	where
	
	\[
	A=
	\begin{bmatrix}
		a_{11} & a_{12} & \cdots & a_{1n}\\
		a_{21} & a_{22} & \cdots & a_{2n}\\
		\vdots & \vdots & \ddots & \vdots\\
		a_{n1} & a_{n2} & \cdots & a_{nn}
	\end{bmatrix}
	\]
	
	\section*{Que 6}
	
	Create the following table with proper alignment, borders, and headings.
	
	\begin{center}
		\begin{tabular}{|c|c|c|c|c|}
			\hline
			\textbf{Roll No.} & \textbf{Name} & \textbf{Department} & \textbf{Semester} & \textbf{CGPA}\\
			\hline
			101 & Amit & CSE & 2 & 8.45\\
			\hline
			102 & Neha & CSE & 2 & 9.12\\
			\hline
			103 & Rahul & ECE & 4 & 8.76\\
			\hline
			104 & Priya & ECE & 4 & 9.01\\
			\hline
		\end{tabular}
	\end{center}
	
	\section*{Que 7}
	
	Create the following table using multicolumn.
	
	\begin{center}
		\begin{tabular}{|c|c|c|}
			\hline
			\multicolumn{3}{|c|}{\textbf{Academic Performance}}\\
			\hline
			Student & Midterm & Final\\
			\hline
			A & 78 & 89\\
			\hline
			B & 85 & 92\\
			\hline
			C & 91 & 94\\
			\hline
		\end{tabular}
	\end{center}
	
	\section*{Que 8}
	
	Create the following table using multirow.
	
	\begin{center}
		\begin{tabular}{|c|l|c|c|}
			\hline
			\textbf{Department} & \textbf{Course} & \textbf{Credits} & \textbf{Grade}\\
			\hline
			\multirow{3}{*}{CSE}
			& Data Structures & 4 & A\\
			\cline{2-4}
			& Machine Learning & 4 & A+\\
			\cline{2-4}
			& Operating Systems & 3 & B+\\
			\hline
			\multirow{2}{*}{ECE}
			& Digital Electronics & 4 & A\\
			\cline{2-4}
			& Signals & 3 & B+\\
			\hline
		\end{tabular}
	\end{center}
	
	\section*{Que 9}
	
	Create the following scientific table using xcolor and colortbl.
	
	\begin{center}
		\begin{tabular}{|l|c|c|c|c|}
			\hline
			\rowcolor{headerblue}
			\color{white}\textbf{Model} &
			\color{white}\textbf{Parameters} &
			\color{white}\textbf{FLOPs} &
			\color{white}\textbf{Accuracy} &
			\color{white}\textbf{Energy (kWh)}\\
			\hline
			\rowcolor{lightblue}
			ResNet18 & 11.7M & 1.8G & 94.2\% & 0.00093\\
			\hline
			MobileNetV2 & 3.5M & 0.3G & 92.8\% & 0.00061\\
			\hline
			\rowcolor{lightblue}
			EfficientNet-B0 & 5.3M & 0.39G & 93.5\% & 0.00072\\
			\hline
		\end{tabular}
	\end{center}
	
	\section*{Que 10}
	
	Create a student-performance table with conditional cell coloring.
	
	\begin{center}
		\begin{tabular}{|c|c|c|c|}
			\hline
			\rowcolor{headerblue}
			\color{white}\textbf{Student} &
			\color{white}\textbf{Mathematics} &
			\color{white}\textbf{Physics} &
			\color{white}\textbf{Computer Science}\\
			\hline
			Amit & \cellcolor{greenmark}92 & \cellcolor{yellowmark}81 & \cellcolor{greenmark}95\\
			\hline
			Neha & \cellcolor{redmark}72 & \cellcolor{yellowmark}88 & \cellcolor{greenmark}91\\
			\hline
			Rahul & \cellcolor{redmark}68 & \cellcolor{redmark}74 & \cellcolor{yellowmark}79\\
			\hline
			Priya & \cellcolor{greenmark}96 & \cellcolor{greenmark}94 & \cellcolor{greenmark}98\\
			\hline
		\end{tabular}
	\end{center}
	
	\textbf{Color Rules:}
	
	\begin{itemize}
		\item Marks $\geq 90$: Green
		\item Marks 75--89: Yellow
		\item Marks $<75$: Red
	\end{itemize}
	
	\section*{Que 11}
	
	Create a table containing paragraphs inside cells.
	
	\begin{center}
		\small
		\renewcommand{\arraystretch}{1.5}
		\begin{tabular}{|p{2.4cm}|p{4cm}|p{4cm}|p{4cm}|}
			\hline
			\textbf{Technique} &
			\textbf{Description} &
			\textbf{Advantages} &
			\textbf{Limitations}\\
			\hline
			Pruning &
			Removes unnecessary parameters from a neural network. &
			Reduces model size and computation. &
			May reduce accuracy.\\
			\hline
			Quantization &
			Represents weights using lower precision. &
			Reduces memory and energy consumption. &
			Can introduce numerical error.\\
			\hline
			Knowledge Distillation &
			Transfers knowledge from a large teacher model to a smaller student model. &
			Produces compact models. &
			Requires training a teacher model.\\
			\hline
		\end{tabular}
	\end{center}
	
	\section*{Que 12}
	
	Create a professional table using the booktabs package.
	
	\begin{table}[ht]
		\centering
		\caption{Comparison of Different Models}
		\label{tab:model-comparison}
		\begin{tabular}{lcccc}
			\toprule
			\textbf{Model} &
			\textbf{Parameters} &
			\textbf{FLOPs} &
			\textbf{Accuracy} &
			\textbf{Energy}\\
			\midrule
			ResNet18 & 11.7M & 1.82G & 94.21 & 0.00093\\
			MobileNetV2 & 3.5M & 0.30G & 92.84 & 0.00061\\
			EfficientNet-B0 & 5.3M & 0.39G & 93.52 & 0.00072\\
			\bottomrule
		\end{tabular}
	\end{table}
	
	Table \ref{tab:model-comparison} presents a comparison of different models.
	
	\section*{Que 13}
	
	Create a table where one cell contains another table.
	
	\begin{center}
		\renewcommand{\arraystretch}{1.3}
		\begin{tabular}{|c|c|c|}
			\hline
			\textbf{Model} & \textbf{Architecture} & \textbf{Performance}\\
			\hline
			CNN & Conv $\rightarrow$ ReLU $\rightarrow$ Pool &
			\begin{tabular}{cc}
				Metric & Value\\
				\hline
				Accuracy & 94\%\\
				F1 Score & 93\%
			\end{tabular}\\
			\hline
			ResNet & Conv $\rightarrow$ BN $\rightarrow$ ReLU &
			\begin{tabular}{cc}
				Metric & Value\\
				\hline
				Accuracy & 96\%\\
				F1 Score & 95\%
			\end{tabular}\\
			\hline
			ViT & Attention $\rightarrow$ MLP &
			\begin{tabular}{cc}
				Metric & Value\\
				\hline
				Accuracy & 97\%\\
				F1 Score & 96\%
			\end{tabular}\\
			\hline
		\end{tabular}
	\end{center}
	
	\section*{Que 14}
	
	Create the following table containing mathematical expressions.
	
	\begin{center}
		\renewcommand{\arraystretch}{1.7}
		\begin{tabular}{|c|c|c|c|}
			\hline
			\textbf{Model} &
			\textbf{Loss Function} &
			\textbf{Parameters} &
			\textbf{Complexity}\\
			\hline
			Linear &
			$L=\dfrac{1}{n}\sum_{i=1}^{n}(y_i-\hat{y}_i)^2$ &
			10K &
			$O(n)$\\
			\hline
			CNN &
			$L=-\sum_i y_i\log(\hat{y})$ &
			1.2M &
			$O(n^2)$\\
			\hline
			Transformer &
			$L=-\sum_i y_i\log p_i$ &
			110M &
			$O(n^2d)$\\
			\hline
		\end{tabular}
	\end{center}
	
	\vspace{1cm}
	
	\begin{center}
		\textbf{--- End of Assignment ---}
	\end{center}
	
\end{document}
